\documentclass[10pt,journal,compsoc]{IEEEtran}

\usepackage[utf8]{inputenc}
\usepackage[T1]{fontenc}
\usepackage[spanish,es-tabla]{babel}
\usepackage{graphicx}
\usepackage{amsmath}
\usepackage{amssymb}
\usepackage{geometry}
\usepackage{hyperref}
\usepackage{listings}
\usepackage{booktabs}
\usepackage{color}

\geometry{margin=1in}

\definecolor{codegreen}{rgb}{0,0.6,0}
\definecolor{codegray}{rgb}{0.5,0.5,0.5}
\definecolor{codepurple}{rgb}{0.58,0,0.82}
\definecolor{backcolour}{rgb}{0.95,0.95,0.92}

\lstdefinestyle{mystyle}{
    backgroundcolor=\color{backcolour},   
    commentstyle=\color{codegreen},
    keywordstyle=\color{magenta},
    numberstyle=\tiny\color{codegray},
    stringstyle=\color{codepurple},
    basicstyle=\ttfamily\footnotesize,
    breakatwhitespace=false,         
    breaklines=true,                 
    captionpos=b,                    
    keepspaces=true,                 
    numbers=left,                    
    numbersep=5pt,                  
    showspaces=false,                
    showstringspaces=false,
    showtabs=false,                  
    tabsize=2
}
\lstset{style=mystyle}

\begin{document}

\title{Manual Técnico de la Arquitectura: Sistema de Verificación de Identidad y Detección de Fraude en Cédulas por WhatsApp}

\author{Esteban Mercado, Juan David Ocampo y Camilo}

\markboth{Manual Técnico del Software,~Vol.~1, No.~1, Julio~2026}%
{Asistente de Detección de Fraude por WhatsApp}

\IEEEtitleabstractindextext{%
\begin{abstract}
Este documento describe formalmente la arquitectura de software, el diseño del grafo de estados conversacional, el modelo de datos persistido en Supabase, la integración con las API de AWS (Lambda y Rekognition), y los modelos de análisis de métricas para la detección de fraude mediante fotos de la cédula colombiana recibidas a través del canal de WhatsApp. El sistema combina técnicas de procesamiento de lenguaje natural generativo (Google Gemini 2.5 Flash) y machine learning de visión artificial (AWS Rekognition) para automatizar el veredicto de autenticidad en tiempo real.
\end{abstract}
}

\maketitle

\IEEEdisplaynontitleabstractindextext
\IEEEpeerreviewmaketitle

\section{Introducción}
\label{sec:intro}

\IEEEpaperstring{La suplantación de identidad y la presentación de documentos falsificados} (tales como fotocopias a color, capturas en pantallas o impresiones de alta fidelidad) constituyen uno de los mayores vectores de fraude en las industrias financiera y de tecnología. El presente software, denominado **NeoAstrum**, propone un enfoque robusto que combina la agilidad y accesibilidad de la aplicación de mensajería **WhatsApp** como interfaz conversacional, con el poder analítico de la computación en la nube.

\subsection{Alcance del Sistema}
El sistema guía al usuario de manera intuitiva para que proporcione secuencialmente las fotos de la parte frontal y trasera de su cédula colombiana. Una vez capturadas, el motor de orquestación coordina la descarga de las imágenes desde los servidores de Meta, las transfiere a un microservicio en AWS Lambda y realiza la validación técnica. 

\subsection{Objetivos del Software}
\begin{itemize}
    \item **Accesibilidad:** Permitir a cualquier ciudadano validar su identidad a través de WhatsApp sin necesidad de instalar aplicaciones adicionales.
    \item **Automatización:** Retornar una clasificación de veracidad (Real vs Fraude/Fotocopia) en menos de 5 segundos.
    \item **Telemetría y Monitoreo:** Registrar cada etapa del flujo con tiempos de ejecución e indicadores clave para su posterior explotación en un panel web interactivo.
\end{itemize}

\section{Arquitectura del Sistema}
\label{sec:architecture}

El sistema NeoAstrum adopta una arquitectura modular distribuida dividida en tres capas principales: la interfaz conversacional, el orquestador de lógica de negocio y los servicios cognitivos de análisis.

\subsection{Capa Orquestadora: FastAPI y LangGraph}
El backend principal está construido utilizando el framework de alto rendimiento **FastAPI**. La gestión de la conversación es controlada por **LangGraph**, que modela el flujo de validación como un grafo de estados dirigido. El estado de la conversación se modela en la estructura de datos `AgentState`:

\begin{lstlisting}[language=Python, caption=Definición de AgentState en el Orquestador]
class AgentState(TypedDict):
    phone: str
    input_text: str | None
    media_id: str | None
    current_state: str  # GREETING, AWAITING_FRONT, AWAITING_BACK, DONE
    front_result: dict | None
    back_result: dict | None
    response: str
    latency_aws_ms: int | None
    latency_gemini_ms: int | None
    node_executed: str | None
    status: str
    is_photocopy: bool
    lambda_score: float | None
    lambda_confidence: float | None
    lambda_response: dict | None
\end{lstlisting}

\subsection{Lógica de Transición de Estados}
El flujo conversacional sigue una máquina de estados determinista guiada por la entrada del usuario:
\begin{enumerate}
    \item **GreetingNode:** Saluda al usuario y solicita la foto frontal.
    \item **ProcessFrontNode:** Descarga la imagen frontal de WhatsApp, invoca la Lambda de AWS y transiciona a `AWAITING_BACK` si la imagen es válida.
    \item **ProcessBackNode:** Descarga la imagen trasera, invoca la Lambda, realiza la consolidación del veredicto final y transiciona a `DONE`.
    \item **ResetNode:** Limpia los datos de sesión y reinicia el flujo cuando el usuario escribe *reiniciar*.
    \item **UnsupportedTextNode:** Responde con asistencia conversacional cuando se reciben mensajes fuera de la secuencia del flujo.
\end{enumerate}

\subsection{Integración con AWS Lambda y Rekognition}
Para la evaluación de las imágenes, el orquestador se conecta de forma segura a una API expuesta en AWS que ejecuta una función Lambda. El flujo técnico de procesamiento es el siguiente:
\begin{enumerate}
    \item **Descarga y Codificación:** El backend descarga los bytes de la imagen de los servidores de Meta, la convierte en formato binario y la transmite mediante una petición `POST` HTTP a la Lambda de AWS.
    \item **AWS Rekognition:** La función Lambda recibe la imagen, invoca las APIs de **AWS Rekognition** para la detección y lectura del documento, y extrae métricas de clasificación de textura y análisis de bordes para detectar anomalías físicas.
    \item **Retorno del Veredicto:** La Lambda responde con un objeto JSON que contiene la predicción, la confianza y el score obtenido.
\end{enumerate}

\subsection{Capa Conversacional: Google Gemini 2.5 Flash}
El orquestador utiliza la API de **Google Gemini 2.5 Flash** para humanizar las respuestas. Gemini recibe el prompt de contexto técnico (el JSON con los resultados de AWS Rekognition) y genera un mensaje amigable y claro en español para el usuario en WhatsApp. 

Para evitar problemas de truncado del texto, se configuró un límite de `max_output_tokens = 1000` en la configuración de generación del modelo.

\section{Modelo de Datos y Persistencia}
\label{sec:database}

El sistema utiliza una base de datos relacional \textbf{PostgreSQL} alojada en \textbf{Supabase} para mantener el estado de las conversaciones y registrar los eventos históricos para su análisis estadístico.

\subsection{Esquema de Tablas}
La base de datos consta de dos tablas principales: \texttt{whatsapp\_sessions} y \texttt{validation\_logs}.

\subsubsection{Tabla: whatsapp\_sessions}
Esta tabla almacena el estado actual de la sesión del usuario. Se indexa por el número de teléfono del cliente para garantizar transiciones de estado rápidas.

\begin{table}[ht]
\centering
\caption{Atributos de la tabla whatsapp\_sessions}
\footnotesize
\renewcommand{\arraystretch}{1.4}
\begin{tabular}{p{1.8cm} p{2cm} p{3cm}}
\toprule
\textbf{Columna} & \textbf{Tipo} & \textbf{Descripción} \\
\midrule
\texttt{phone} & TEXT (PK) & Número de teléfono del usuario. \\
\texttt{state} & TEXT & Estado actual de la máquina de estados. \\
\texttt{front\_result} & JSONB & Metadatos técnicos de la cara frontal. \\
\texttt{back\_result} & JSONB & Metadatos técnicos de la cara trasera. \\
\texttt{updated\_at} & TIMESTAMPTZ & Fecha y hora del último mensaje. \\
\bottomrule
\end{tabular}
\end{table}

\subsubsection{Tabla: validation\_logs}
Esta tabla registra cada una de las interacciones y transacciones procesadas por el bot. Es la fuente de datos principal para el Dashboard de monitoreo.

\begin{table}[ht]
\centering
\caption{Atributos de la tabla validation\_logs}
\footnotesize
\renewcommand{\arraystretch}{1.4}
\begin{tabular}{p{2.2cm} p{1.8cm} p{2.8cm}}
\toprule
\textbf{Columna} & \textbf{Tipo} & \textbf{Descripción} \\
\midrule
\texttt{id} & SERIAL (PK) & Identificador único correlativo. \\
\texttt{phone} & TEXT & Número de teléfono del usuario. \\
\texttt{timestamp} & TIMESTAMP\-TZ & Fecha y hora del procesamiento. \\
\texttt{node} & TEXT & Nombre del nodo ejecutado. \\
\texttt{whatsapp\_\allowbreak message\_id} & TEXT & Identificador wamid de Meta. \\
\texttt{lambda\_\allowbreak score} & NUMERIC & Score de fraude obtenido. \\
\texttt{lambda\_\allowbreak confidence} & NUMERIC & Confianza de detección de cédula. \\
\texttt{lambda\_\allowbreak response} & JSONB & Respuesta JSON completa de AWS. \\
\texttt{latency\_\allowbreak aws\_ms} & INTEGER & Latencia del API de AWS (ms). \\
\texttt{latency\_\allowbreak gemini\_ms} & INTEGER & Latencia del API de Gemini (ms). \\
\texttt{status} & TEXT & Estado del resultado de la petición. \\
\texttt{is\_photocopy} & BOOLEAN & Indicador binario de fraude. \\
\bottomrule
\end{tabular}
\end{table}

\subsection{Restricciones de Pooling en Supabase}
Dado que Supabase implementa un Connection Pooler (PgBouncer) configurado en \textbf{modo transacción} para manejar altas concurrencias (puerto \texttt{6543}), se configuró explícitamente el parámetro \texttt{statement\_cache\_size=0} en la librería \texttt{asyncpg} en el backend. 

Esta configuración es obligatoria para evitar errores de ejecución en la base de datos debido al reuso indebido de prepared statements a nivel de sesión en el pooler.

\section{Modelo de Métricas y Fórmulas Matemáticas}
\label{sec:metrics}

Las métricas mostradas en el Dashboard de Detection Attacks se derivan de la consolidación de los registros históricos en la tabla `validation\_logs`. Esta sección detalla la construcción conceptual y matemática de dichos indicadores.

\subsection{Criterio de Clasificación de Fraude}
El sistema asigna un veredicto binario sobre la veracidad del documento basándose en la puntuación de confianza $S$ retornada por el microservicio de AWS Rekognition. El modelo está calibrado para que $S \in [0, 1]$, donde un valor cercano a $0$ representa un documento real y un valor cercano a $1$ indica una sospecha o intento de fraude.

Definimos el veredicto de fraude como una función indicadora $I_{\text{fraude}}$ con un umbral de decisión $\tau = 0.5$:

\begin{equation}
I_{\text{fraude}}(S) = \begin{cases} 
1 & \text{si } S \ge 0.5 \quad \text{(FRAUDE / Sospecha)} \\
0 & \text{si } S < 0.5 \quad \text{(REAL / Aprobado)}
\end{cases}
\end{equation}

Este umbral coincide con la línea de decisión visible en los histogramas de distribución del dashboard.

\subsection{Indicadores Clave de Rendimiento (KPIs)}
Sea $N_{\text{total}}$ el número total de transacciones registradas en el webhook, y sea $M$ el subconjunto de transacciones correspondientes al análisis de imágenes frontales ($N_{\text{front}}$) o traseras ($N_{\text{back}}$):

\begin{equation}
M = \{x \in \text{logs} \mid \text{node}(x) \in \{\text{'process\_front'}, \text{'process\_back'}\} \}
\end{equation}

Definimos los indicadores de volumen del siguiente modo:
\begin{align}
\text{Total Validations} &= N_{\text{total}} \\
\text{Predicted Real} &= \sum_{i \in M} (1 - I_{\text{fraude}}(S_i)) \\
\text{Predicted Fraud} &= \sum_{i \in M} I_{\text{fraude}}(S_i)
\end{align}

Las tasas porcentuales de clasificación se formulan como:
\begin{equation}
\text{Percentage of Real (\%)} = \left( \frac{\text{Predicted Real}}{|M|} \right) \times 100
\end{equation}

\begin{equation}
\text{Percentage of Fraud (\%)} = \left( \frac{\text{Predicted Fraud}}{|M|} \right) \times 100
\end{equation}

\subsection{Histogramas de Distribución de Scores}
Para analizar la distribución probabilística del clasificador, se agrupan los scores continuos de cada transacción en $K=10$ bins de ancho constante $W = 0.1$. El conteo de registros $C_k$ en el bin $k$ (donde $k \in \{1, 2, \dots, 10\}$) se define como:

\begin{equation}
C_k = |\{i \in M \mid S_i \in [ (k-1) \cdot 0.1, \ k \cdot 0.1 ) \}|
\end{equation}
*(Para el último bin $k=10$, el intervalo es cerrado por la derecha $[0.9, 1.0]$).*

\subsection{Promedios de Latencia y Tiempos de Respuesta}
Sea $L_{\text{AWS}, i}$ la latencia en milisegundos de la transacción $i$ en AWS Lambda, y sea $L_{\text{Gemini}, i}$ la latencia del servicio de humanización de Gemini. Los promedios aritméticos generales mostrados en el dashboard se calculan como:

\begin{equation}
\bar{L}_{\text{AWS}} = \frac{1}{|M|} \sum_{i \in M} L_{\text{AWS}, i}
\end{equation}

\begin{equation}
\bar{L}_{\text{Gemini}} = \frac{1}{|M|} \sum_{i \in M} L_{\text{Gemini}, i}
\end{equation}


\end{document}
